% ChargeShield-FL --- DSN 2027 --- 3. System, Data and Threat Model
% Scritta il 2026-09-23. Fonti: docs/SISTEMA.md (sez. 1-4, 7-9), STATO.md
% sez. 3.7 per gli identificativi (corregge l'inversione Caltech/JPL di
% setup.tex, segnalazione 36), fonti giornalistiche/vendor per gli incidenti
% (verificate il 2026-09-23, vedi references.bib).

\section{System, Data and Threat Model}
\label{sec:system}

\subsection{Incidents in the EV-Charging Ecosystem}
\label{sec:threat-incidents}

Table~\ref{tab:incidents} summarises four publicly reported incidents. We
include them because they establish two premises of this work, and we are
explicit about a third point that they do \emph{not} establish.

\begin{table*}[t]
\centering
\caption{Publicly reported incidents involving EV-charging data. Figures are
those reported by the cited sources and were not verified independently.}
\label{tab:incidents}
\footnotesize
\begin{tabular}{@{}p{1.3cm}p{2.6cm}p{3.0cm}p{5.2cm}p{4.0cm}@{}}
\toprule
When & Affected party & Vector & Data reported exposed & Relation to this study \\
\midrule
Jun.\ 2023 & Shell Recharge Solutions~\cite{whittaker2023shell} & Internet-facing logging database without password, found by a researcher & fleet customers' names, e-mail addresses and phone numbers; VINs; locations of charging points, including residential ones & central store of charging logs; outside our threat model \\
Nov.\ 2024 & several charge-point operators using a common third-party charging application~\cite{robb2024cpo} & data offered on a hacking forum & about 116{,}000 records: names, contact details, addresses, vehicle make and model, VINs, keys and tokens, precise geolocation, OCPP logs; reportedly payment information & central store; shows session and protocol logs linked to identity \\
Mar.\ 2026 & ELECQ, charger manufacturer~\cite{page2026elecq} & ransomware on the cloud platform; databases encrypted and copied & names, e-mail addresses, phone numbers, home addresses; the company stated that no financial data was involved and that chargers were unaffected; notified to the UK ICO and the German BfDI & central cloud store; outside our threat model \\
Recurrent (reports since 2024) & drivers at public chargers~\cite{muncaster2024quishing,rac2024qr} & fake QR codes on chargers leading to look-alike payment pages (``quishing'') & payment card details and, potentially, credentials & targets the payment interface, not training; outside our threat model \\
\bottomrule
\end{tabular}
\end{table*}

First, charging data sit next to identity. The records exposed in 2023 and
2024 joined sessions, locations and protocol logs to names, addresses and
VINs; this is what makes the membership of a single session a personal fact
rather than a statistical curiosity. Second, centralised stores of such data
fail in practice, through misconfiguration (2023), compromise of a shared
application (2024) or ransomware (2026), which is the usual motivation for
keeping records at the operator that collected them and training federatively.

What the incidents do not establish is equally important. None of them is a
membership inference attack, and none involves a machine-learning model: they
compromise a database, a cloud platform or the user. FL does not remove these
vectors --- each operator still holds its own records and payment interfaces
--- and DP, which bounds what a trained model or update reveals about a unit,
protects neither a database nor a payment page. Quishing, in particular, lies
entirely outside the training pipeline and is listed only to delimit the
scope. The question we study is narrower: what the artifacts \emph{added} by
federated training can reveal about the records behind them.

\subsection{Data}
\label{sec:system-data}

We use ACN-Data~\cite{lee2019acndata} from its three real sites, all available
years, for 66{,}713 sessions (Table~\ref{tab:data}); each session carries its
site, and each site is one FL client. Sessions are split 80/20 into training
and hold-out at random with the run's seed; the training part of a site is
that client's data. Six continuous features describe a session: delivered
energy, mean power (energy over duration), requested energy, declared
availability window, connection hour in the site's local time, and duration.
They are scaled to $[0,1]$ with min--max statistics computed on the training
split only. ACN-Data carries no intrusion labels (the field provided for them
is empty in every session), so we cannot and do not evaluate the detector's
ability to detect attacks.

\begin{table}[t]
\centering
\caption{Data per site. Identifier coverage: share of sessions carrying a user
identifier (measured on 2026-09-21).}
\label{tab:data}
\footnotesize
\begin{tabular}{@{}lrr@{}}
\toprule
Site (client) & Sessions & With user identifier \\
\midrule
Caltech  & 31{,}404 & 52.1\% \\
JPL      & 33{,}629 & 93.5\% \\
Office~1 & 1{,}680  & 34.8\% \\
\midrule
Total    & 66{,}713 & 72.6\% \\
\bottomrule
\end{tabular}
\end{table}

Among identified sessions there are 1{,}028 distinct users, with a median of
14 sessions each, and 8.5\% of them charge at more than one site. The median
is not a maximum, and coverage differs across sites, so these figures
characterise the data rather than evaluate the privacy of persons
(Section~\ref{sec:threat-unit}).

\subsection{Federation and Model}
\label{sec:system-model}

The model is a dense autoencoder $6 \rightarrow 16 \rightarrow 8 \rightarrow 4
\rightarrow 8 \rightarrow 16 \rightarrow 6$ with 570 linear parameters and 48
normalisation parameters, batch normalisation in the encoder (group
normalisation under record-level DP, which is incompatible with batch
statistics), a sigmoid output and mean-squared-error loss. It flags a session
as anomalous when its reconstruction error is high. We chose an autoencoder
because no labels exist, because the nodes that train it are station
controllers, and because its reconstruction error is at once the detector's
score and the quantity that loss-based attacks threshold, which removes a
degree of freedom between what the system computes and what the attacker
observes.

Training uses Adam with learning rate $10^{-3}$ and batch size 32. Each round
every client runs 50 local epochs with the FedProx objective, $\mu =
0.01$~\cite{li2020fedprox}, the proximal term being active from the second
round; the server averages the updates weighted by sample count. The ordinary
protocol has 10 rounds with all three clients in every round. The canary
regime (Section~\ref{sec:method-controls}) changes these settings on purpose:
3 rounds, 1{,}000 local epochs, hidden widths $(32,16)$ with latent dimension
8, and a seventh feature, the session start time, which is nearly unique per
session.

\subsection{Differential Privacy Mechanisms}
\label{sec:system-dp}

\paragraph{Client-level} At the end of a round, the client's update --- the
difference from the weights it received --- is clipped to $L_2$ norm $C = 1$
and perturbed with Gaussian noise of standard deviation
$\sigma = C\sqrt{2\ln(1.25/\delta)}/\varepsilon$, with $\delta = 10^{-5}$;
batch-normalisation buffers are excluded from clipping and noise. The
mechanism can be placed at three points, which determine what the aggregator
receives (Table~\ref{tab:placements}). The per-round $\varepsilon$ of this
mechanism is a calibration parameter for $\sigma$ rather than a verified
guarantee, for three reasons: the classical Gaussian calibration above
guarantees $(\varepsilon,\delta)$-DP only for $\varepsilon <
1$~\cite[Thm.~A.1]{dwork2014algorithmic}, whereas the sweep uses $\varepsilon$
up to 16 (an analytic calibration~\cite{balle2018improving} would be needed to
state the guarantee there); the batch-normalisation buffers are released
without noise; and the normalisation statistics are computed on training data
outside the accounted mechanism. We therefore report $\sigma$ together with
$\varepsilon$, and cumulative values $T\varepsilon$ only as labels.

\begin{table}[t]
\centering
\caption{Client-level DP placements and what the aggregator receives.}
\label{tab:placements}
\footnotesize
\begin{tabular}{@{}p{1.45cm}p{1.55cm}p{1.9cm}p{2.25cm}@{}}
\toprule
Placement & Clips & Adds noise & Aggregator receives \\
\midrule
none & --- & --- & A0: raw update, no DP \\
\texttt{dp-fedavg} & server, per client, on arrival & server, per client & A1: raw update \\
\texttt{central} & client & server, once on the aggregate & A2: clipped, not noised \\
\texttt{local} & client & client & A3: clipped and noised \\
\bottomrule
\end{tabular}
\end{table}

\paragraph{Record-level} DP-SGD with microbatches of one record, per-example
gradient clipping and Gaussian noise $\sigma C$ on the sum, the proximal term
being added after the noise. It protects the unit that the attacks measure.
Its budget is computed with an R\'enyi accountant~\cite{mironov2017renyi}
under a Poisson-sampling assumption; the current accountant under-counts the
rounds, so record-level $\varepsilon$ values are not reported until it is
corrected \Susy{segnalazione 4: fl\_rounds e delta da cfg["experiment"], n per
client, Poisson contro shuffle}. With record-level DP the three observation
points collapse into one. The two mechanisms are enabled one at a time.

\subsection{Threat Model}
\label{sec:threat-model}

We evaluate one scenario: an \emph{honest-but-curious} aggregator that follows
the protocol and records what it legitimately receives, together with any
participant that receives the global model. The two correspond to two
observation surfaces. The \emph{global model} is available to every
participant; an attack on it can place a record in the training data of
\emph{some} client. The \emph{per-client update} is available to the
aggregator; an attack on it attributes membership to a specific operator.
Figure~\ref{fig:observation} shows the four observation points of the update.

% ChargeShield-FL --- DSN 2027 --- Figura: punti di osservazione A0-A3
% 2026-09-23. Sostituisce nel nuovo main.tex figures/adversary_observation_points.tex
% (che resta nel progetto): quella didascalia diceva che l'update grezzo non
% lascia mai il client, mentre sotto dp-fedavg il server lo riceve e lo
% clippa all'arrivo (SISTEMA.md sez. 4, tabella dei placement).
% Nessun \\ annidato dentro un altro comando nei nodi (bug TikZ gia' visto).
\begin{figure}[t]
\centering
\begin{tikzpicture}[
  font=\footnotesize,
  stage/.style={draw, rounded corners, minimum height=0.65cm,
                minimum width=1.45cm, align=center, font=\scriptsize},
  adv/.style={draw, circle, minimum size=0.5cm, thick,
              red!65!black, font=\scriptsize\bfseries},
  arr/.style={-Latex, thick}
]
  \node[stage] (raw)   {update\\$g$};
  \node[stage, right=0.5cm of raw]   (clip)  {clipped\\$\bar{g}$};
  \node[stage, right=0.5cm of clip]  (noise) {noised\\$\tilde{g}$};
  \node[stage, right=0.5cm of noise] (agg)   {aggregate\\FedAvg};

  \draw[arr] (raw)   -- node[above, font=\scriptsize]{clip} (clip);
  \draw[arr] (clip)  -- node[above, font=\scriptsize]{$+\mathcal{N}$} (noise);
  \draw[arr] (noise) -- (agg);

  \node[adv, above=0.35cm of raw]   (a1) {A1};
  \node[adv, above=0.35cm of clip]  (a2) {A2};
  \node[adv, above=0.35cm of noise] (a3) {A3};
  \draw[arr, red!65!black] (a1) -- (raw);
  \draw[arr, red!65!black] (a2) -- (clip);
  \draw[arr, red!65!black] (a3) -- (noise);

  \node[font=\scriptsize, align=center, below=0.3cm of raw]   {\texttt{dp-fedavg}\\(A0 if no DP)};
  \node[font=\scriptsize, align=center, below=0.3cm of clip]  {\texttt{central}};
  \node[font=\scriptsize, align=center, below=0.3cm of noise] {\texttt{local}};
\end{tikzpicture}
\caption{Observation points of a client's update along the client-level
privatisation pipeline, ordered by the information left to the adversary.
The three placements are the same Gaussian mechanism observed at different
points, not three mechanisms. Under record-level DP the noise is added inside
local training and the points collapse into one.}
\label{fig:observation}
\end{figure}


A0 is the raw update of a run without DP. A1 is the raw update under
\texttt{dp-fedavg}, which the server receives before clipping and noising it:
relative to the \texttt{central} and \texttt{local} deployments it is a
stronger position, since noise added after the observation cannot protect
against that observation. A2 is the clipped update that a server adding noise
centrally legitimately receives. A3 is the clipped and noised update, which is
also what an observer on the network would see.
\Susy{SISTEMA.md sez. 4 e 8 descrive A1 una volta come ``cio' che vede
l'aggregatore'' e una volta come ``insider sul client'': ho usato la prima
lettura, coerente con la tabella dei placement. Da confermare.}

\paragraph{Adversary knowledge} The adversary knows the architecture and the
training procedure. The shadow models of LiRA are trained on the site's real
data, members and hold-out included, and the shadow attack trains its single
shadow on half of the members. Both are broader accesses than the threat model
grants to an aggregator, so these attacks measure a generous upper bound on
what that aggregator could infer, not an operational attack.

\paragraph{Out of scope} A malicious client (an integrity threat), a
compromised gateway, a network eavesdropper under transport encryption and an
adaptive attacker tailored to this system are not evaluated, nor are the
vectors of Section~\ref{sec:threat-incidents}. The original project also asked
whether the attack could be detected from anomalies in the gradients. We
exclude that question: an honest-but-curious aggregator runs the inference
offline without altering any update, so a detector that sees only the updates
has no signal that an inference took place, and similarity-based filters such
as Krum~\cite{blanchard2017machine} are not detectors of passive membership
inference.

\subsection{Protected Unit and Person}
\label{sec:threat-unit}

Neither protected unit coincides with the person. Record-level DP bounds each
session separately, but a user contributes a median of 14 sessions, whose
guarantees compose. Client-level DP bounds a site's whole contribution, but a
user who charges at several sites contributes to several clients, and no
operator can clip data it does not hold. The dependence among a person's
records also weakens the membership bound of
Section~\ref{sec:bg-dp}~\cite{humphries2020differentially}. We measure these
properties of the data; we do not evaluate subject-level protection or
subject-level attacks.
